\documentclass[a4paper, serif, 10pt]{moderncv}

%% moderncv
% style and color
\moderncvstyle{banking}
\moderncvcolor{black}

%% page layout/margins
\usepackage[top=2.5cm, bottom=2.5cm, left=1.5cm, right=1.5cm]{geometry}

\nopagenumbers{}

%% Babel config for German language
% ref:
% - https://latex3.github.io/babel/guides/locale-german.html
\usepackage[ngerman]{babel}

%% fontspec
\usepackage{fontspec}

\IfFontExistsTF{Linux Libertine}{
  \setmainfont[Ligatures={TeX, Common}, Numbers=OldStyle]{Linux Libertine}
}{}
\IfFontExistsTF{Linux Libertine O}{
  \setmainfont[Ligatures={TeX, Common}, Numbers=OldStyle]{Linux Libertine O}
}{}

%% CTeX
% CJK support, used to show CJK characters in the resume
%
% - fontset=none: disable builtin fontset but instead set the CJK font manually
% - heading=false: disable ctex heading
% - punct=kaiming: use kaiming punctuations styles for CJK
% - scheme=plain: use plain scheme, do not override \normalsize font size
% - space=auto: space settings for CJK characters
%
% ref:
% - http://ctan.mirrorcatalogs.com/language/chinese/ctex/ctex.pdf
\usepackage[UTF8, heading=false, punct=kaiming, scheme=plain, space=auto]{ctex}

\IfFontExistsTF{Noto Serif CJK SC}{
  \setCJKmainfont{Noto Serif CJK SC}
}{}
\IfFontExistsTF{Noto Sans CJK SC}{
  \setCJKsansfont{Noto Sans CJK SC}
}{}

%% line spacing
\usepackage{setspace}
\setstretch{1.125}

%% URL styling - use normal text instead of monospace
\urlstyle{same}

% Auto-underline all links
% Defer until \begin{document} so hyperref is already loaded
\AtBeginDocument{%
  \let\oldhref\href
  \renewcommand{\href}[2]{\oldhref{#1}{\underline{#2}}}%
}

%% Patch \httplink and \httpslink to handle full URLs
% The original moderncv macros always prepend http:// or https:// to URLs.
% This causes issues when the URL already has a protocol (e.g., https://example.com)
% resulting in malformed URLs like https://https://example.com.
% This patch checks if the URL already contains "://" and uses it directly if so.
% Note: We use \str_set:Nx (x-expansion) to expand macros like \@homepage first.
\ExplSyntaxOn
\str_new:N \g_yamlresume_url_str

\RenewDocumentCommand{\httpslink}{O{}m}{
  \str_set:Nx \g_yamlresume_url_str {#2}
  \str_if_in:NnTF \g_yamlresume_url_str {://}
    { \url{#2} }
    { \url{https://#2} }
}

\RenewDocumentCommand{\httplink}{O{}m}{
  \str_set:Nx \g_yamlresume_url_str {#2}
  \str_if_in:NnTF \g_yamlresume_url_str {://}
    { \url{#2} }
    { \url{http://#2} }
}
\ExplSyntaxOff

\name{Dr. Anna Schmidt}{}
\title{Senior Projektmanagerin | Agile Transformation \& Digitalisierung}
\phone[mobile]{+49 30 9876543}
\email{anna.schmidt@example.de}
\homepage{https://www.annaschmidt-pm.de}

\address{Friedrichstraße 123, Berlin, Berlin, Deutschland, 10117}{}{}

\extrainfo{{\small \faLinkedin}\ \href{https://www.linkedin.com/in/annaschmidtpm}{@annaschmidtpm} {} {} {} • {} {} {} 
{\small \faGithub}\ \href{https://github.com/annaschmidt}{@annaschmidt} {} {} {} • {} {} {} 
{\small \faTwitter}\ \href{https://twitter.com/annaschmidt\_pm}{@annaschmidt\_pm}}

\begin{document}

\maketitle

\section{Basisinformationen}

\cvline{}{Senior Projektmanagerin mit über 8 Jahren Erfahrung in der Leitung komplexer IT- und Digitalisierungsprojekte. Nachweisliche Erfolge in der agilen Transformation, im Stakeholder-Management und in der Budgetverantwortung von bis zu 5 Mio. EUR.
%
\begin{itemize}
\item Führung interdisziplinärer Teams von bis zu 25 Mitarbeitenden
\item Zertifizierte Scrum Masterin und PMP mit fundierten Kenntnissen in SAFe und Kanban
\item Starke Kommunikationsfähigkeiten in Deutsch und Englisch, verhandlungssicher in Französisch
\item Leidenschaft für nachhaltige Prozessoptimierung und datengetriebene Entscheidungen
\end{itemize}}
\section{Ausbildung}

\cventry{Okt. 2013–Sept. 2015}
        {Master, Wirtschaftsinformatik, Punkte: 1.7}
        {Technische Universität München}
        {\url{https://www.tum.de}}
        {}
        {\begin{itemize}
\item Schwerpunkt auf IT-Projektmanagement und agile Methoden
\item Masterarbeit: "Erfolgsfaktoren agiler Transformationsprojekte in mittelständischen Unternehmen"
\end{itemize}
\textbf{Kurse}: Projektmanagement, Agile Softwareentwicklung, Datenbanken, Unternehmensführung, IT-Controlling}

\cventry{Okt. 2010–Sept. 2013}
        {Bachelor, Betriebswirtschaftslehre, Punkte: 1.9}
        {Universität Mannheim}
        {\url{https://www.uni-mannheim.de}}
        {}
        {\begin{itemize}
\item Grundlagen in BWL, VWL und Informatik
\item Auslandssemester an der Universität St. Gallen
\end{itemize}
\textbf{Kurse}: Betriebswirtschaftslehre, Statistik, Informatik}
\section{Berufserfahrung}

\cventry{Jan. 2021–Heute}
        {Senior Projektmanagerin}
        {Digital Pioneers GmbH}
        {\url{https://www.digital-pioneers.de}}
        {}
        {\begin{itemize}
\item Leitung eines portfolioweiten agilen Transformationsprogramms mit 12 Projekten und einem Budget von 5 Mio. EUR
\item Einführung eines skalierbaren SAFe-Frameworks, Steigerung der Liefergeschwindigkeit um 30\%
\item Coaching von 5 Scrum-Teams und Aufbau eines internen PMO
\item Berichterstattung an den Vorstand und Steuerung externer Dienstleister
\end{itemize}
\textbf{Schlüsselwörter}: Agile Transformation, SAFe, PMO, Budgetverantwortung, Stakeholder-Management}

\cventry{Juli 2017–Dez. 2020}
        {Projektmanager}
        {TechSolutions AG}
        {\url{https://www.techsolutions.de}}
        {}
        {\begin{itemize}
\item Erfolgreiche Leitung von 15+ IT-Projekten in den Bereichen Cloud-Migration und SAP-Einführung
\item Verantwortung für Projektplanung, Risikomanagement und Qualitätssicherung
\item Enge Zusammenarbeit mit Product Ownern und Entwicklungsteams
\item Einführung von Kanban in der Softwareentwicklung, Reduzierung der Durchlaufzeit um 25\%
\end{itemize}
\textbf{Schlüsselwörter}: Cloud-Migration, SAP, Kanban, Risikomanagement, Jira}

\cventry{Aug. 2015–Juni 2017}
        {Junior Projektmanager}
        {StartUp Innovations GmbH}
        {\url{https://www.startup-innovations.de}}
        {}
        {\begin{itemize}
\item Unterstützung bei der Planung und Umsetzung von Digitalisierungsprojekten
\item Erstellung von Projektplänen, Statusberichten und Präsentationen
\item Koordination von Meetings und Workshops
\end{itemize}
\textbf{Schlüsselwörter}: Digitalisierung, Projektplanung, Reporting}
\section{Sprachen}

\cvline{Deutsch}{Muttersprache oder zweisprachig \hfill \textbf{Schlüsselwörter}: Muttersprache}
\cvline{Englisch}{Verhandlungssichere Kenntnisse \hfill \textbf{Schlüsselwörter}: Verhandlungssicher, TOEFL 110}
\cvline{Französisch}{Gute Kenntnisse \hfill \textbf{Schlüsselwörter}: Grundkenntnisse}
\section{Fähigkeiten}

\cvline{Projektmanagement}{Experte \hfill \textbf{Schlüsselwörter}: Agiles Projektmanagement, Wasserfall, Hybrid, PMI, PRINCE2}
\cvline{Agile Methoden}{Experte \hfill \textbf{Schlüsselwörter}: Scrum, Kanban, SAFe, Lean, OKR}
\cvline{Tools}{Erfahren \hfill \textbf{Schlüsselwörter}: Jira, Confluence, MS Project, Trello, Miro}
\cvline{Führung \& Kommunikation}{Erfahren \hfill \textbf{Schlüsselwörter}: Teamentwicklung, Moderation, Konfliktmanagement, Stakeholder-Management}
\section{Auszeichnungen}

\cventry{Dez. 2023}
        {Digital Pioneers GmbH}
        {Projektmanagerin des Jahres 2023}
        {}
        {}
        {Auszeichnung für die herausragende Leitung des agilen Transformationsprogramms und die erfolgreiche Einführung des PMO.}

\cventry{Okt. 2015}
        {Technische Universität München}
        {Best Master Thesis Award}
        {}
        {}
        {Auszeichnung für die beste Masterarbeit im Bereich Wirtschaftsinformatik.}
\section{Zertifikate}

\cventry{März 2020}
        {Project Management Institute}
        {Project Management Professional (PMP)}
        {\url{https://www.pmi.org/certifications/project-management-pmp}}
        {}
        {}

\cventry{Juni 2019}
        {Scrum.org}
        {Professional Scrum Master I}
        {\url{https://www.scrum.org/professional-scrum-master-i-certification}}
        {}
        {}

\cventry{Sept. 2021}
        {Scaled Agile}
        {SAFe 5 Agilist}
        {\url{https://www.scaledagile.com/certification/safe-agilist/}}
        {}
        {}
\section{Veröffentlichungen}

\cventry{Mai 2018}
        {Erfolgsfaktoren agiler Transformationsprojekte}
        {Projektmanagement aktuell}
        {\url{https://www.pm-aktuell.de/artikel/erfolgsfaktoren-agiler-transformation}}
        {}
        {\begin{itemize}
\item Analyse von 20 agilen Transformationsprojekten in mittelständischen Unternehmen
\item Identifikation von kritischen Erfolgsfaktoren wie Führung, Kommunikation und Kultur
\item Handlungsempfehlungen für Projektmanager und Führungskräfte
\end{itemize}}
\section{Projekte}

\cventry{März 2021–Nov. 2022}
        {Aufbau eines zentralen PMO zur Steuerung von 12 agilen Projekten}
        {Einführung eines agilen PMO}
        {\url{https://www.digital-pioneers.de/projekte/agiles-pmo}}
        {}
        {\begin{itemize}
\item Etablierung von standardisierten Prozessen und Reporting
\item Einführung von Jira und Confluence als zentrale Plattformen
\item Schulung von 50 Mitarbeitenden in agilen Methoden
\item Reduzierung der Projektdurchlaufzeit um 20\%
\end{itemize}
\textbf{Schlüsselwörter}: PMO, Agile, Jira, Confluence}

\cventry{Jan. 2019–Dez. 2020}
        {Migration von 200+ Anwendungen in die AWS-Cloud}
        {Cloud-Migration für einen DAX-Konzern}
        {\url{https://www.techsolutions.de/projekte/cloud-migration}}
        {}
        {\begin{itemize}
\item Leitung eines Teams von 15 Entwicklern und Architekten
\item Sicherstellung der Termin- und Budgettreue
\item Erfolgreiche Abnahme durch den Kunden
\end{itemize}
\textbf{Schlüsselwörter}: AWS, Cloud, Migration}
\section{Interessen}

\cvline{Sport}{Marathon, Skifahren, Wandern}
\cvline{Musik}{Klavier, Gitarre}
\section{Ehrenamtliche Tätigkeiten}

\cventry{Jan. 2020–Dez. 2023}
        {Mentorin}
        {PMI Germany Chapter}
        {\url{https://www.pmi-germany.de}}
        {}
        {\begin{itemize}
\item Mentoring von Nachwuchsprojektmanagern
\item Organisation von Netzwerkveranstaltungen und Workshops
\item Vorträge zu agilen Methoden und Karriereentwicklung
\end{itemize}}
    
\end{document}